
\documentclass[aps,preprint]{revtex4}
%%%%%%%%%%%%%%%%%%%%%%%%%%%%%%%%%%%%%%%%%%%%%%%%%%%%%%%%%%%%%%%%%%%%%%%%%%%%%%%%%%%%%%%%%%%%%%%%%%%%%%%%%%%%%%%%%%%%%%%%%%%%%%%%%%%%%%%%%%%%%%%%%%%%%%%%%%%%%%%%%%%%%%%%%%%%%%%%%%%%%%%%%%%%%%%%%%%%%%%%%%%%%%%%%%%%%%%%%%%%%%%%%%%%%%%%%%%%%%%%%%%%%%%%%%%%
\usepackage{amsfonts}
\usepackage{amsmath}
\usepackage{amssymb}
\usepackage{graphicx}

\setcounter{MaxMatrixCols}{10}
%TCIDATA{OutputFilter=LATEX.DLL}
%TCIDATA{Version=5.50.0.2960}
%TCIDATA{<META NAME="SaveForMode" CONTENT="1">}
%TCIDATA{BibliographyScheme=Manual}
%TCIDATA{Created=Friday, April 06, 2012 11:15:25}
%TCIDATA{LastRevised=Thursday, September 10, 2026 07:32:36}
%TCIDATA{<META NAME="GraphicsSave" CONTENT="32">}
%TCIDATA{<META NAME="DocumentShell" CONTENT="Articles\SW\REVTeX 4">}
%TCIDATA{Language=American English}
%TCIDATA{CSTFile=revtex4.cst}

\newtheorem{theorem}{Theorem}
\newtheorem{acknowledgement}[theorem]{Acknowledgement}
\newtheorem{algorithm}[theorem]{Algorithm}
\newtheorem{axiom}[theorem]{Axiom}
\newtheorem{claim}[theorem]{Claim}
\newtheorem{conclusion}[theorem]{Conclusion}
\newtheorem{condition}[theorem]{Condition}
\newtheorem{conjecture}[theorem]{Conjecture}
\newtheorem{corollary}[theorem]{Corollary}
\newtheorem{criterion}[theorem]{Criterion}
\newtheorem{definition}[theorem]{Definition}
\newtheorem{example}[theorem]{Example}
\newtheorem{exercise}[theorem]{Exercise}
\newtheorem{lemma}[theorem]{Lemma}
\newtheorem{notation}[theorem]{Notation}
\newtheorem{problem}[theorem]{Problem}
\newtheorem{proposition}[theorem]{Proposition}
\newtheorem{remark}[theorem]{Remark}
\newtheorem{solution}[theorem]{Solution}
\newtheorem{summary}[theorem]{Summary}
\newenvironment{proof}[1][Proof]{\noindent\textbf{#1.} }{\ \rule{0.5em}{0.5em}}
\input{tcilatex}
\begin{document}

\title{}
\author{}
\maketitle

\begin{abstract}
We characterize the necessary and sufficient conditions for a partial normal
state defined on the subspace $\mathcal{A}_4 = \mathcal{A}^0 \oplus \mathcal{%
A}^2 \oplus \mathcal{A}^4 \subset C\ell_{2r}^+(\mathbb{C})$ (grades 0, 2,
and 4 of the even subalgebra of a complex Clifford algebra) to extend to a
full state on the entire even subalgebra. Using the explicit coefficient
expansion in terms of Clifford generators and the complete characterization
of grade-4 inhomogeneous projectors, we derive a linear programming
condition. The introduction of a number operator $N$ with integrality
condition $\mathrm{Tr}(Nd) = n \in \mathbb{Z}_{>0}$ forces the partial state
into the interior of the state space, guaranteeing multiple extensions,
including the canonical grand canonical ensemble $D_\beta = e^{-\beta N}/%
\mathrm{Tr}(e^{-\beta N})$.
\end{abstract}

\tableofcontents

\section{Introduction}

\section{Preliminaries}

\subsection{Complex Clifford Algebras}

\begin{enumerate}
\item Let $\mathcal{F}=C\ell _{2r}(\mathbb{C})$ be the complex Clifford
algebra over $\mathbb{C}^{2r}$ with generators $\{x_{1},\ldots ,x_{2r}\}$

\begin{enumerate}
\item satisfying:

\begin{enumerate}
\item 
\begin{equation}
x_{i}x_{j}+x_{j}x_{i}=2\delta _{ij}I_{\mathcal{F}}.
\end{equation}%
which gives

\item 
\begin{equation}
x_{i}x_{j}=-x_{j}x_{i};\quad 1\leq i\neq j\leq 2r
\end{equation}

\item 
\begin{equation}
x_{j}^{2}=I_{\mathcal{F}};\quad 1\leq j\leq 2r
\end{equation}
\end{enumerate}

\item It is isomorphic to 
\begin{equation}
\mathcal{F}\cong \mathrm{Mat}(2^{r},\mathbb{C}).
\end{equation}
\end{enumerate}

\item $\mathcal{F}$ can be viewed as a CAR algebra generated by the

\begin{enumerate}
\item Majorano operators%
\begin{equation}
\{x_{1},\ldots ,x_{2r}\}
\end{equation}%
via the operator $\equiv $ geometric algebra product.

\item Creation operators 
\begin{equation}
\{a_{1}^{\dagger },\ldots ,a_{r}^{\dagger }\}
\end{equation}%
via the operator product and the adjoint operation "$\dagger "$.

\item Ket-Bra operators 
\begin{equation}
\{\left\vert \varphi _{1}\right\rangle \left\langle \phi \right\vert ,\ldots
,\left\vert \varphi _{r}\right\rangle \left\langle \phi \right\vert \}
\end{equation}%
via the Greu$\beta $ exterior algebra product "$\bullet "$the adjoint
operation $"\dagger ".$
\end{enumerate}

\item These bases generate the direct sum decompositions

\begin{enumerate}
\item \textbf{Graded Clifford algebra}

\begin{enumerate}
\item 
\begin{equation}
\mathcal{F}=\bigoplus\limits_{0\leq n\leq 2r}\mathcal{F}_{n},
\end{equation}%
where 
\begin{equation}
\mathcal{F}_{n}==\mathbb{C}\text{-}\func{linspan}\left\{ x_{1},\ldots
,x_{n}\right\} ,
\end{equation}%
thus using conventional vector space notation 
\begin{equation}
\mathcal{F}_{1}=\mathcal{F}^{1},
\end{equation}%
which gives in standard exterior algebra notation

\item 
\begin{equation}
\mathcal{F}_{n}={\Huge \wedge }^{n}\mathcal{F}_{1}=\mathcal{F}^{n};\quad
0\leq n\leq 2r
\end{equation}%
leading to

\item 
\begin{equation}
\mathcal{F}=\bigoplus\limits_{0\leq n\leq 2r}\mathcal{F}^{n}=\bigwedge
\left( \mathcal{F}^{1}\right) .
\end{equation}

\item The grade 1 elements 
\begin{equation}
\mathcal{F}^{1}=\mathcal{F}_{1}=\mathbb{C}\text{-}\func{linspan}%
\{x_{1},\ldots ,x_{2r}\}
\end{equation}%
form a complex vector space thus one can define 
\begin{equation}
\mathcal{F}^{n}={\Huge \wedge }^{n}\mathcal{F}^{1}
\end{equation}%
and observe that 
\begin{equation}
\mathcal{F}=\bigoplus\limits_{0\leq n\leq 2r}\mathcal{F}^{n}=\bigwedge
\left( \mathcal{F}^{1}\right) .
\end{equation}
\end{enumerate}

\item \textbf{Graded CAR algebra}

\begin{enumerate}
\item In terms of creation and annihilation operators 
\begin{equation}
\mathcal{F}=\bigoplus\limits_{0\leq m,n\leq r}\mathcal{A}_{mn}=\bigoplus%
\limits_{0\leq p\leq r}\mathcal{A}_{p},
\end{equation}%
where 
\begin{equation}
\mathcal{A}_{mn}=\mathbb{C}\text{-}\func{linspan}\left\{ \left.
a_{j_{1}}^{\dagger }\cdots a_{j_{m}}^{\dagger }a_{k_{1}}\cdots
a_{k_{n}}\right\vert 1\leq j_{1}<\cdots <j_{m}\leq r,1\leq k_{1}<\cdots
<k_{n}\leq r\right\}
\end{equation}%
and 
\begin{equation}
\mathcal{A}_{p}=\mathbb{C}\text{-}\func{linspan}\left\{ \left.
a_{j_{1}}^{\dagger }\cdots a_{j_{m}}^{\dagger }a_{j_{1}}\cdots
a_{j_{n}}\right\vert 1\leq j_{1}<\cdots <j_{m}\leq r,1\leq k_{1}<\cdots
<k_{n}\leq r,0\leq m,n\leq r,m+n=p\right\} .
\end{equation}

\item The grade 1 elements 
\begin{equation}
\mathcal{A}^{1}=\mathcal{A}_{1}=\mathbb{C}\text{-}\func{linspan}%
\{a_{1}^{\dagger },\ldots ,a_{r}^{\dagger },a_{1},\ldots ,a_{r}\}
\end{equation}%
form a complex vector space thus one can define 
\begin{equation}
\mathcal{A}^{n}={\Huge \wedge }^{n}\mathcal{A}^{1}
\end{equation}%
and observe that 
\begin{equation}
\mathcal{F}=\bigoplus\limits_{0\leq n\leq 2r}\mathcal{A}^{n}=\bigwedge
\left( \mathcal{A}^{1}\right) .
\end{equation}

\begin{remark}
Note that the exponent takes value all the way to $2r.$
\end{remark}

\begin{example}
\begin{equation}
\mathcal{A}_{2}=\mathbb{C}\text{-}\func{linspan}\left\{ \left.
a_{j_{1}}^{\dagger }a_{j_{m}}^{\dagger }a_{j_{1}}a_{j_{n}}\right\vert 1\leq
j_{1}<\cdots <j_{m}\leq r,1\leq k_{1}<\cdots <k_{n}\leq r,0\leq m,n\leq
r,m+n=2\right\}
\end{equation}%
and 
\begin{equation}
\mathcal{A}_{2}=\mathcal{A}_{20}\oplus \mathcal{A}_{11}\oplus \mathcal{A}%
_{02}.
\end{equation}
\end{example}

\begin{example}
\begin{equation}
\mathcal{A}^{2}=\mathbb{C}\text{-}\func{linspan}\left\{ \left. 
\begin{array}{c}
a_{j_{1}}^{\dagger }\wedge a_{j_{2}}^{\dagger }\wedge a_{j_{3}}^{\dagger
}\wedge a_{j_{4}}^{\dagger }=a_{j_{1}}^{\dagger }a_{j_{2}}^{\dagger
}a_{j_{3}}^{\dagger }a_{j_{4}}^{\dagger },\,1\leq
j_{1}<j_{2}<j_{3}<j_{4}\leq r, \\ 
a_{k_{1}}\wedge a_{k_{2}}\wedge a_{k_{3}}\wedge
a_{k_{4}}=a_{k_{1}}a_{k_{2}}a_{k_{3}}a_{k_{4}},1\leq
k_{1}<k_{2}<k_{3}<k_{4}\leq r, \\ 
a_{j_{1}}^{\dagger }\wedge a_{j_{2}}^{\dagger }\wedge a_{j_{3}}^{\dagger
}\wedge a_{k_{1}}=a_{j_{1}}^{\dagger }a_{j_{2}}^{\dagger }a_{j_{3}}^{\dagger
}\wedge a_{k_{1}},1\leq j_{1}<j_{2}<j_{3}\leq r,\,1\leq k_{1}\leq r, \\ 
a_{j_{1}}^{\dagger }\wedge a_{j_{2}}^{\dagger }\wedge a_{k_{1}}\wedge
a_{k_{2}}=a_{j_{1}}^{\dagger }a_{j_{2}}^{\dagger }\wedge
a_{k_{1}}a_{k_{2}},1\leq j_{1}<j_{2}<j_{3}\leq r,\,1\leq k_{1}<k_{2}\leq r,
\\ 
a_{j_{1}}^{\dagger }\wedge a_{k_{1}}\wedge a_{k_{2}}\wedge
a_{k_{3}}=a_{j_{1}}^{\dagger }\wedge a_{k_{1}}a_{k_{2}}a_{k_{3}},1\leq
j_{1}\leq r,\,1\leq k_{1}<k_{2}<k_{3}\leq r%
\end{array}%
\right\vert \right\}
\end{equation}
\end{example}

\item \textbf{NOTE\ THAT }%
\begin{equation}
\mathcal{A}^{n}\neq \mathcal{A}_{n},\quad 2\leq n\leq 2r,
\end{equation}%
which can be seen by 
\begin{eqnarray}
a_{j}^{\dagger }\wedge a_{j} &=&\frac{1}{2}\left( a_{j}^{\dagger
}a_{j}-a_{j}a_{j}^{\dagger }\right) =\frac{1}{2}\left( a_{j}^{\dagger
}a_{j}-I_{\mathcal{F}}+a_{j}^{\dagger }a_{j}\right) \\
&=&a_{j}^{\dagger }a_{j}-\frac{1}{2}I_{\mathcal{F}}
\end{eqnarray}%
and equivalently 
\begin{equation}
a_{j}^{\dagger }a_{j}=a_{j}^{\dagger }\wedge a_{j}+\frac{1}{2}I_{\mathcal{F}%
}\in \mathcal{A}_{even4}^{\left( 2\right) }=\mathcal{A}^{0}\oplus \mathcal{A}%
^{2}.
\end{equation}%
This produces 
\begin{equation}
a_{j}^{\dagger }a_{j}a_{k}^{\dagger }a_{k}=\left( a_{j}^{\dagger }\wedge
a_{j}+\frac{1}{2}I_{\mathcal{F}}\right) \left( a_{k}^{\dagger }\wedge a_{k}+%
\frac{1}{2}I_{\mathcal{F}}\right) =a_{j}^{\dagger }\wedge a_{j}\wedge
a_{k}^{\dagger }\wedge a_{k}+\frac{1}{2}a_{j}^{\dagger }\wedge a_{j}+\frac{1%
}{2}a_{k}^{\dagger }\wedge a_{k}+\frac{1}{4}I_{\mathcal{F}}
\end{equation}%
and hence 
\begin{equation}
a_{j}^{\dagger }a_{k}^{\dagger }a_{k}a_{j}=-\left( a_{j}^{\dagger }\wedge
a_{k}^{\dagger }\wedge a_{k}\wedge a_{j}+\frac{1}{2}a_{j}^{\dagger }\wedge
a_{j}+\frac{1}{2}a_{k}^{\dagger }\wedge a_{k}+\frac{1}{4}I_{\mathcal{F}%
}\right) \in \mathcal{A}_{even}^{\left( 4\right) }=\mathcal{A}^{0}\oplus 
\mathcal{A}^{2}\oplus \mathcal{A}^{4}.
\end{equation}

\item consider $j_{1}=k_{1},j_{2}\neq k_{2}$%
\begin{equation}
a_{j_{1}}^{\dagger }\wedge a_{j_{2}}^{\dagger }\wedge a_{j_{1}}\wedge
a_{k_{2}}=-a_{j_{1}}^{\dagger }\wedge a_{j_{1}}\wedge a_{j_{2}}^{\dagger
}\wedge a_{k_{2}}=-\left( a_{j_{1}}^{\dagger }a_{j_{1}}-\frac{1}{2}I_{%
\mathcal{F}}\right) a_{j_{2}}^{\dagger }a_{k_{2}}=-a_{j_{1}}^{\dagger
}a_{j_{1}}a_{j_{2}}^{\dagger }a_{k_{2}}+\frac{1}{2}a_{j_{2}}^{\dagger
}a_{k_{2}}
\end{equation}%
and $j_{1}=k_{1},j_{2}=k_{2}$ 
\begin{eqnarray}
a_{j_{1}}^{\dagger }\wedge a_{j_{2}}^{\dagger }\wedge a_{j_{1}}\wedge
a_{j_{2}} &=&-a_{j_{1}}^{\dagger }\wedge a_{j_{1}}\wedge a_{j_{2}}^{\dagger
}\wedge a_{j_{2}}=-\left( a_{j_{1}}^{\dagger }a_{j_{1}}-\frac{1}{2}I_{%
\mathcal{F}}\right) \left( a_{j_{2}}^{\dagger }a_{j_{2}}-\frac{1}{2}I_{%
\mathcal{F}}\right) \\
&=&-a_{j_{1}}^{\dagger }a_{j_{1}}a_{j_{2}}^{\dagger
}a_{j_{2}}+a_{j_{1}}^{\dagger }a_{j_{1}}+a_{j_{2}}^{\dagger }a_{j_{2}}-\frac{%
1}{4}I_{\mathcal{F}}
\end{eqnarray}%
showing that 
\begin{equation}
\mathcal{A}^{4}={\Huge \wedge }^{4}\mathcal{A}^{1}\subset \mathcal{A}%
_{1}\oplus \mathcal{A}_{2}\oplus \mathcal{A}_{4}.
\end{equation}

\item Recall%
\begin{equation}
\mathcal{F}=\bigoplus\limits_{0\leq n\leq 2r}\mathcal{A}^{n}={\Huge \wedge }%
\left( \mathcal{A}^{1}\right) =\mathcal{A}^{\left( 2r\right) },
\end{equation}%
where 
\begin{equation}
\mathcal{A}^{\left( n\right) }=\bigoplus\limits_{0\leq p\leq n}\mathcal{A}%
^{p}.
\end{equation}
\end{enumerate}

\item \textbf{Graded Ket-Bra algebra}%
\begin{equation}
\mathcal{F}=\bigoplus\limits_{0\leq m,n\leq r}\mathcal{B}\left( \mathcal{H}%
^{m},\mathcal{H}^{n}\right) ,
\end{equation}%
where 
\begin{equation}
\mathcal{B}\left( \mathcal{H}^{m},\mathcal{H}^{n}\right) =\mathbb{C}\text{-}%
\func{linspan}\left\{ \left. \left\vert \varphi _{j_{1}}\right\rangle
\left\langle \phi \right\vert \bullet \ldots \bullet \left\vert \varphi
_{j_{m}}\right\rangle \left\langle \phi \right\vert \bullet \left\vert \phi
\right\rangle \left\langle \varphi _{k_{1}}\right\vert \bullet \ldots
\bullet \left\vert \phi \right\rangle \left\langle \varphi
_{k_{n}}\right\vert \right\vert 1\leq j_{1}<\cdots <j_{m}\leq r,1\leq
k_{1}<\cdots <k_{n}\leq r\right\} .
\end{equation}
\end{enumerate}

\item \textbf{Relationships between the different decompositions:}

\begin{enumerate}
\item Equality of subspaces in the CAR algebra and Clifford algebra 
\begin{equation}
\mathcal{F}_{n}=\mathcal{F}^{n}=\mathcal{A}^{n},\quad 0\leq n\leq 2r
\end{equation}

\item Bra-Kets in terms of second quantized operators 
\begin{eqnarray}
\left\vert \varphi _{j_{1}}\right\rangle \left\langle \phi \right\vert
\bullet \cdots \bullet \left\vert \varphi _{j_{m}}\right\rangle \left\langle
\phi \right\vert \bullet \left\vert \phi \right\rangle \left\langle \varphi
_{k_{1}}\right\vert \bullet \cdots \bullet \left\vert \phi \right\rangle
\left\langle \varphi _{k_{n}}\right\vert &=&\left\vert \varphi
_{j_{1}}\wedge \cdots \wedge \varphi _{j_{m}}\right\rangle \left\langle
\varphi _{k_{1}}\wedge \cdots \wedge \varphi _{k_{n}}\right\vert \\
&=&\sqrt{m!n!}\left\vert \varphi _{j_{1}}\cdots \varphi
_{j_{m}}\right\rangle \left\langle \varphi _{k_{1}}\cdots \varphi
_{k_{n}}\right\vert =P_{\mathcal{H}^{m}}a_{j_{1}}^{\dagger }\cdots
a_{j_{m}}^{\dagger }a_{j_{1}}\cdots a_{j_{n}}P_{\mathcal{H}^{n}}
\end{eqnarray}%
i.e. 
\begin{equation}
\mathcal{B}\left( \mathcal{H}^{m},\mathcal{H}^{n}\right) =P_{\mathcal{H}^{m}}%
\mathcal{A}P_{\mathcal{H}^{n}}=\widehat{P_{\mathcal{H}^{m}\mathcal{H}^{n}}}%
\left( \mathcal{A}\right) =\widehat{P_{\mathcal{H}^{m}\mathcal{H}^{n}}}%
\left( \mathcal{A}_{mn}\right) ;\quad 0\leq m,n\leq r.
\end{equation}
\end{enumerate}

\begin{example}
\begin{enumerate}
\item 
\begin{eqnarray}
\left\vert \varphi _{j_{1}}\cdots \varphi _{j_{m}}\right\rangle \left\langle
\varphi _{k_{1}}\cdots \varphi _{k_{n}}\right\vert  &=&P_{\mathcal{H}%
^{m}}a_{j_{1}}^{\dagger }\cdots a_{j_{m}}^{\dagger }a_{k_{1}}\cdots
a_{k_{n}}P_{\mathcal{H}^{n}} \\
&=&\left\vert \varphi _{j_{1}}\cdots \varphi _{j_{m}}\right\rangle
\left\langle \left. \varphi _{j_{1}}\cdots \varphi _{j_{m}}\right\vert
a_{j_{1}}^{\dagger }\cdots a_{j_{m}}^{\dagger }a_{j_{1}}\cdots
a_{j_{n}}\varphi _{k_{1}}\cdots \varphi _{k_{n^{{}}}}\right\rangle
\left\langle \varphi _{k_{1}}\cdots \varphi _{k_{n}}\right\vert 
\end{eqnarray}

\item 
\begin{equation}
\left\vert \varphi _{j_{1}}\cdots \varphi _{j_{m^{\prime }}}\right\rangle
\left\langle \varphi _{k_{1}}\cdots \varphi _{k_{n^{\prime }}}\right\vert
=P_{\mathcal{H}^{m^{\prime }}}a_{j_{1}}^{\dagger }\cdots a_{j_{m}}^{\dagger
}a_{k_{1}}\cdots a_{k_{n}}P_{\mathcal{H}^{n^{\prime }}}=\left\vert \varphi
_{j_{1}}\cdots \varphi _{j_{m^{\prime }}}\right\rangle \left\langle \left.
\varphi _{j_{1}}\cdots \varphi _{j_{m^{\prime }}}\right\vert
a_{j_{1}}^{\dagger }\cdots a_{j_{m}}^{\dagger }a_{j_{1}}\cdots
a_{j_{n}}\varphi _{k_{1}}\cdots \varphi _{k_{n^{\prime }}}\right\rangle
\left\langle \varphi _{k_{1}}\cdots \varphi _{k_{n^{\prime }}}\right\vert \ 
\end{equation}

\item Using a general expansion of the projector 
\begin{equation}
P_{\mathcal{H}^{m}}=\sum_{1\leq j\leq \binom{r}{m}}\left\vert \Psi
_{mj}\right\rangle \left\langle \Psi _{mj}\right\vert
\end{equation}%
one obtains 
\begin{eqnarray}
\widehat{P_{\mathcal{H}^{m}\mathcal{H}^{n}}}\left( a_{j_{1}}^{\dagger
}\cdots a_{j_{m}}^{\dagger }a_{k_{1}}\cdots a_{k_{n}}\right) &=&P_{\mathcal{H%
}^{m}}a_{j_{1}}^{\dagger }\cdots a_{j_{m}}^{\dagger }a_{k_{1}}\cdots
a_{k_{n}}P_{\mathcal{H}^{n}}=\sum_{1\leq j\leq \binom{r}{m}}\left\vert \Psi
_{mj}\right\rangle \left\langle \Psi _{mj}\right\vert \left(
a_{j_{1}}^{\dagger }\cdots a_{j_{m}}^{\dagger }a_{k_{1}}\cdots
a_{k_{n}}\right) \sum_{1\leq k\leq \binom{r}{n}}\left\vert \Psi
_{mk}\right\rangle \left\langle \Psi _{mk}\right\vert \\
&=&\sum_{\substack{ 1\leq j\leq \binom{r}{m}  \\ 1\leq k\leq \binom{r}{n}}}%
\left\vert \Psi _{mj}\right\rangle \left\langle \Psi _{nk}\right\vert
\left\langle \left. \Psi _{mj}\right\vert a_{j_{1}}^{\dagger }\cdots
a_{j_{m}}^{\dagger }a_{k_{1}}\cdots a_{k_{n}}\Psi _{nk}\right\rangle
\end{eqnarray}%
and hence%
\begin{eqnarray}
\widehat{P_{\mathcal{H}^{m^{\prime }}\mathcal{H}^{n^{\prime }}}}\left(
a_{j_{1}}^{\dagger }\cdots a_{j_{m}}^{\dagger }a_{k_{1}}\cdots
a_{k_{n}}\right) P_{\mathcal{H}^{n^{\prime }}} &=&P_{\mathcal{H}^{m^{\prime
}}}a_{j_{1}}^{\dagger }\cdots a_{j_{m}}^{\dagger }a_{k_{1}}\cdots
a_{k_{n}}P_{\mathcal{H}^{n^{\prime }}}=\sum_{\substack{ 1\leq j\leq \binom{r%
}{m}  \\ 1\leq k\leq \binom{r}{m}}}\left\vert \Psi _{m^{\prime
}j}\right\rangle \left\langle \Psi _{n^{\prime }k}\right\vert \left\langle
\left. \Psi _{m^{\prime }j}\right\vert a_{j_{1}}^{\dagger }\cdots
a_{j_{m}}^{\dagger }a_{k_{1}}\cdots a_{k_{n}}\Psi _{n^{\prime
}k}\right\rangle \\
&=&\sum_{\substack{ 1\leq j\leq \binom{r}{m}  \\ 1\leq k\leq \binom{r}{m}}}%
\left\vert \Psi _{m^{\prime }j}\right\rangle \left\langle \Psi _{n^{\prime
}k}\right\vert L_{m^{\prime }n^{\prime }}^{mn}\left( \left\vert \Psi
_{m^{\prime }j}\right\rangle \left\langle \Psi _{n^{\prime }k}\right\vert
\right) _{j_{1}\cdots j_{m}k_{1}\cdots k_{m}}.
\end{eqnarray}%
\textbf{where the Reduced Transition matrix elements are given by} 
\begin{equation}
L_{m^{\prime }n^{\prime }}^{mn}\left( \left\vert \Psi _{m^{\prime
}j}\right\rangle \left\langle \Psi _{n^{\prime }k}\right\vert \right)
_{j_{1}\cdots j_{m}k_{1}\cdots k_{m}}=\left\langle \left. \Psi _{m^{\prime
}j}\right\vert a_{j_{1}}^{\dagger }\cdots a_{j_{m}}^{\dagger
}a_{k_{1}}\cdots a_{k_{n}}\Psi _{n^{\prime }k}\right\rangle .
\end{equation}

\item Let $m=n$

\begin{enumerate}
\item 
\begin{equation}
\left\langle \left. \Psi _{mj}\right\vert a_{j_{1}}^{\dagger }\cdots
a_{j_{m}}^{\dagger }a_{k_{1}}\cdots a_{k_{m}}\Psi _{mk}\right\rangle
=L_{m}^{m}\left( \left\vert \Psi _{mj}\right\rangle \left\langle \Psi
_{mk}\right\vert \right) _{j_{1}\cdots j_{m}jk_{1}\cdots k_{m}}
\end{equation}

\item Let $m=n$ $\neq m^{\prime }=n^{\prime }$ 
\begin{eqnarray}
\left\langle \left. \Psi _{m^{\prime }j}\right\vert a_{j_{1}}^{\dagger
}\cdots a_{j_{m}}^{\dagger }a_{k_{1}}\cdots a_{k_{m}}\Psi _{m^{\prime
}k}\right\rangle &=&L_{m^{\prime }}^{m}\left( \left\vert \Psi _{m^{\prime
}j}\right\rangle \left\langle \Psi _{m^{\prime }k}\right\vert \right)
_{j_{1}\cdots j_{m}jk_{1}\cdots k_{m}},\quad 0\leq m^{\prime }\leq m \\
\left\langle \left. \Psi _{m^{\prime }j}\right\vert a_{j_{1}}^{\dagger
}\cdots a_{j_{m}}^{\dagger }a_{k_{1}}\cdots a_{k_{m}}\Psi _{m^{\prime
}k}\right\rangle &=&L_{m^{\prime }}^{m}\left( \left\vert \Psi _{m^{\prime
}j}\right\rangle \left\langle \Psi _{m^{\prime }k}\right\vert \right)
_{j_{1}\cdots j_{m}jk_{1}\cdots k_{m}}=0,\quad m^{\prime }+1\leq m\leq r
\end{eqnarray}

\item In general 
\begin{eqnarray}
\left\langle \left. \Psi _{m^{\prime }j}\right\vert a_{j_{1}}^{\dagger
}\cdots a_{j_{m}}^{\dagger }a_{k_{1}}\cdots a_{k_{n}}\Psi _{n^{\prime
}k}\right\rangle &=&L_{m^{\prime }n^{\prime }}^{mn}\left( \left\vert \Psi
_{m^{\prime }j}\right\rangle \left\langle \Psi _{m^{\prime }k}\right\vert
\right) _{j_{1}\cdots j_{m}jk_{1}\cdots k_{m}},\quad 0\leq m^{\prime }\leq
m;0\leq n^{\prime }\leq n \\
\left\langle \left. \Psi _{m^{\prime }j}\right\vert a_{j_{1}}^{\dagger
}\cdots a_{j_{m}}^{\dagger }a_{k_{1}}\cdots a_{k_{m}}\Psi _{m^{\prime
}k}\right\rangle &=&L_{m^{\prime }n^{\prime }}^{mn}\left( \left\vert \Psi
_{m^{\prime }j}\right\rangle \left\langle \Psi _{m^{\prime }k}\right\vert
\right) _{j_{1}\cdots j_{m}jk_{1}\cdots k_{m}}=0,\quad m^{\prime }+1\leq
m\leq r,0\leq n^{\prime }\leq n
\end{eqnarray}

\item \textbf{The reduced }$\left( m,n\right) $-\textbf{order} \textbf{%
transition densities }$D_{m^{\prime }n\prime }^{mn}\left( jk\right) $\textbf{%
\ are defined as} 
\begin{equation*}
D_{m^{\prime }n\prime }^{mn}\left( jk\right) =L_{m^{\prime }n^{\prime
}}^{mn}\left( \left\vert \Psi _{m^{\prime }j}\right\rangle \left\langle \Psi
_{n^{\prime }k}\right\vert \right) ,\quad 1\leq j\leq \binom{m^{\prime }}{m}%
,1\leq k\leq \binom{n^{\prime }}{n},\,0\leq m\leq m^{\prime }\leq r,\,0\leq
n\leq n^{\prime }\leq r\,.
\end{equation*}
\end{enumerate}
\end{enumerate}
\end{example}

\item The \textbf{partial isometry transition operator }$\left\vert \varphi
_{j_{1}}\cdots \varphi _{j_{m^{\prime }}}\right\rangle \left\langle \varphi
_{k_{1}}\cdots \varphi _{k_{n^{\prime }}}\right\vert $ as \textbf{%
projections of second quantized operators and of the generalized gxpansion
map}

\begin{enumerate}
\item 
\begin{eqnarray}
\left\vert \varphi _{j_{1}}\cdots \varphi _{j_{m^{\prime }}}\right\rangle
\left\langle \varphi _{k_{1}}\cdots \varphi _{k_{n^{\prime }}}\right\vert
&=&P_{\mathcal{H}^{m^{\prime }}}a_{j_{1}}^{\dagger }\cdots
a_{j_{m}}^{\dagger }a_{k_{1}}\cdots a_{k_{n}}P_{\mathcal{H}^{n^{\prime }}}=%
\widehat{P_{\mathcal{H}^{m^{\prime }}\mathcal{H}^{n^{\prime }}}}\left(
a_{j_{1}}^{\dagger }\cdots a_{j_{m}}^{\dagger }a_{k_{1}}\cdots
a_{k_{n}}\right) \\
&=&\left\vert \varphi _{j_{1}}\cdots \varphi _{j_{m^{\prime }}}\right\rangle
\left\langle \varphi _{k_{1}}\cdots \varphi _{k_{n^{\prime }}}\right\vert
=P_{\mathcal{H}^{m^{\prime }}}\Gamma \left( \left\vert \varphi
_{j_{1}}\cdots \varphi _{j_{m}}\right\rangle \left\langle \varphi
_{k_{1}}\cdots \varphi _{k_{n}}\right\vert \right) P_{\mathcal{H}^{n^{\prime
}}}=\widehat{P_{\mathcal{H}^{m^{\prime }}\mathcal{H}^{n^{\prime }}}}\left(
\Gamma \left( \left\vert \varphi _{j_{1}}\cdots \varphi
_{j_{m}}\right\rangle \left\langle \varphi _{k_{1}}\cdots \varphi
_{k_{n}}\right\vert \right) \right) .
\end{eqnarray}

\item 
\begin{equation}
\Gamma \left( \left\vert \varphi _{j_{1}}\cdots \varphi
_{j_{m}}\right\rangle \left\langle \varphi _{k_{1}}\cdots \varphi
_{k_{n}}\right\vert \right) =a_{j_{1}}^{\dagger }\cdots a_{j_{m}}^{\dagger
}a_{k_{1}}\cdots a_{k_{n}}.
\end{equation}

\item 
\begin{equation}
1\leq j\leq \binom{m^{\prime }}{m},1\leq k\leq \binom{n^{\prime }}{n}%
,\,0\leq m\leq m^{\prime }\leq r,\,0\leq n\leq n^{\prime }\leq r.
\end{equation}

\item The generalized expansion map is invertibable.
\end{enumerate}

\begin{example}
\begin{enumerate}
\item 
\begin{enumerate}
\item giving%
\begin{equation}
\left\vert \varphi _{j_{1}}\cdots \varphi _{j_{m}}\right\rangle \left\langle
\varphi _{k_{1}}\cdots \varphi _{k_{n}}\right\vert =\Gamma ^{-1}\left(
a_{j_{1}}^{\dagger }\cdots a_{j_{m}}^{\dagger }a_{k_{1}}\cdots
a_{k_{n}}\right) .
\end{equation}

\item and \textbf{notably} 
\begin{equation}
\frame{$\left\langle \left. \Psi _{m^{\prime }j}\right\vert \Gamma
^{-1}\left( \left\vert \varphi _{j_{1}}\cdots \varphi _{j_{m^{\prime
}}}\right\rangle \left\langle \varphi _{k_{1}}\cdots \varphi _{k_{n^{\prime
}}}\right\vert \right) \Psi _{n^{\prime }k}\right\rangle =L_{m^{\prime
}n^{\prime }}^{mn}\left( \left\vert \Psi _{m^{\prime }j}\right\rangle
\left\langle \Psi _{n^{\prime }k}\right\vert \right) _{j_{1}\cdots
j_{m}jk_{1}\cdots k_{m}}$}
\end{equation}
\end{enumerate}
\end{enumerate}
\end{example}

\item \textbf{Relationship between matrix elements of reduced transition
operators and matrix elements of second quantized operators} 
\begin{gather}
\left\langle \left. \Psi _{m^{\prime }j}\right\vert a_{j_{1}}^{\dagger
}\cdots a_{j_{m}}^{\dagger }a_{k_{1}}\cdots a_{k_{n}}\Psi _{n^{\prime
}k}\right\rangle =L_{m^{\prime }n^{\prime }}^{mn}\left( \left\vert \Psi
_{m^{\prime }j}\right\rangle \left\langle \Psi _{n^{\prime }k}\right\vert
\right) _{j_{1}\cdots j_{m}jk_{1}\cdots k_{m}} \\
0\leq m^{\prime }\leq m;0\leq n^{\prime }\leq n.
\end{gather}

\item \textbf{Ket-Bra operators in terms of SQ and Majorano operators} 
\begin{gather}
\left\vert \varphi _{j_{1}}\cdots \varphi _{j_{p}}\right\rangle \left\langle
\varphi _{k_{1}}\cdots \varphi _{k_{q}}\right\vert \\
\parallel \\
a_{j_{1}}^{\dagger }\cdots a_{j_{q}}^{\dagger }\left( \dprod\limits_{1\leq
j\leq r}h_{j}\right) a_{k_{p}}\cdots a_{k_{1}} \\
\parallel \\
\dprod\limits_{1\leq \kappa \leq q}a_{j_{\kappa }}^{\dagger
}\dprod\limits_{1\leq j\leq r}h_{j}\dprod\limits_{p\leq \kappa \leq
1}a_{k_{\kappa }} \\
\parallel \\
\dprod\limits_{1\leq \kappa \leq q}\left( \frac{1}{\sqrt{2}}\left( x_{\kappa
_{q}}-ix_{\kappa _{q}+r}\right) \right) \dprod\limits_{1\leq j\leq r}\left( 
\frac{1}{2}I_{\mathcal{F}}-ix_{j}x_{j+r}\right) \dprod\limits_{p\leq \kappa
\leq 1}\left( \frac{1}{\sqrt{2}}\left( x_{k_{\kappa }}+ix_{k_{\kappa
}+r}\right) \right) ,
\end{gather}%
thus \emph{the most general primitive projector }is given by 
\begin{align}
& \left\vert \Psi \right\rangle \left\langle \Psi \right\vert \\
& =\sum_{0\leq p,q\leq r}\sum_{\substack{ 1\leq j_{1}<\cdots <j_{q}\leq r 
\\ 1\leq k_{1}<\cdots <k_{p}\leq r}}\bar{c}_{k_{1}\cdots
k_{p}}c_{j_{1}\cdots j_{q}}\left\vert \varphi _{j_{1}}\cdots \varphi
_{j_{p}}\right\rangle \left\langle \varphi _{k_{1}}\cdots \varphi
_{k_{q}}\right\vert \\
& =\sum_{0\leq p,q\leq r}\sum_{\substack{ 1\leq j_{1}<\cdots <j_{q}\leq r 
\\ 1\leq k_{1}<\cdots <k_{p}\leq r}}\bar{c}_{k_{1}\cdots
k_{p}}c_{j_{1}\cdots j_{q}}\dprod\limits_{1\leq \kappa \leq q}a_{j_{\kappa
}}^{\dagger }\dprod\limits_{1\leq j\leq r}h_{j}\dprod\limits_{p\leq \kappa
\leq 1}a_{k_{\kappa }} \\
& =\sum_{0\leq p,q\leq r}\sum_{\substack{ 1\leq j_{1}<\cdots <j_{q}\leq r 
\\ 1\leq k_{1}<\cdots <k_{p}\leq r}}\bar{c}_{k_{1}\cdots
k_{p}}c_{j_{1}\cdots j_{q}} \\
& \qquad \times \dprod\limits_{1\leq \kappa \leq q}\left( \frac{1}{\sqrt{2}}%
\left( x_{\kappa _{q}}-ix_{\kappa _{q}+r}\right) \right)
\dprod\limits_{1\leq j\leq r}\left( \frac{1}{2}I_{\mathcal{F}%
}-ix_{j}x_{j+r}\right) \dprod\limits_{p\leq \kappa \leq 1}\left( \frac{1}{%
\sqrt{2}}\left( x_{k_{\kappa }}+ix_{k_{\kappa }+r}\right) \right) .
\end{align}

\begin{enumerate}
\item 
\begin{equation}
\left\vert \varphi _{j_{1}}\cdots \varphi _{j_{p}}\right\rangle \left\langle
\varphi _{k_{1}}\cdots \varphi _{k_{q}}\right\vert =\dprod\limits_{1\leq
\kappa \leq q}a_{j_{\kappa }}^{\dagger }\dprod\limits_{1\leq j\leq
r}h_{j}\dprod\limits_{p\leq \kappa \leq 1}a_{k_{\kappa }}=P_{\mathcal{H}%
^{p}}\dprod\limits_{1\leq \kappa \leq q}a_{j_{\kappa }}^{\dagger
}\dprod\limits_{p\leq \kappa \leq 1}a_{k_{\kappa }}P_{\mathcal{H}^{q}}=%
\widehat{P_{\mathcal{H}^{p}\mathcal{H}^{q}}}\left( \dprod\limits_{1\leq
\kappa \leq q}a_{j_{\kappa }}^{\dagger }\dprod\limits_{p\leq \kappa \leq
1}a_{k_{\kappa }}\right)
\end{equation}

\item 
\begin{equation}
\dprod\limits_{1\leq \kappa \leq q}a_{j_{\kappa }}^{\dagger
}\dprod\limits_{p\leq \kappa \leq 1}a_{k_{\kappa }}=\Gamma \left( \left\vert
\varphi _{j_{1}}\cdots \varphi _{j_{p}}\right\rangle \left\langle \varphi
_{k_{1}}\cdots \varphi _{k_{q}}\right\vert \right)
\end{equation}

\item 
\begin{equation}
\Gamma ^{-1}\left( \dprod\limits_{1\leq \kappa \leq q}a_{j_{\kappa
}}^{\dagger }\dprod\limits_{p\leq \kappa \leq 1}a_{k_{\kappa }}\right)
=\left\vert \varphi _{j_{1}}\cdots \varphi _{j_{p}}\right\rangle
\left\langle \varphi _{k_{1}}\cdots \varphi _{k_{q}}\right\vert =P_{\mathcal{%
H}^{p}}\dprod\limits_{1\leq \kappa \leq q}a_{j_{\kappa }}^{\dagger
}\dprod\limits_{p\leq \kappa \leq 1}a_{k_{\kappa }}P_{\mathcal{H}^{q}}=%
\widehat{P_{\mathcal{H}^{p}\mathcal{H}^{q}}}\left( \dprod\limits_{1\leq
\kappa \leq q}a_{j_{\kappa }}^{\dagger }\dprod\limits_{p\leq \kappa \leq
1}a_{k_{\kappa }}\right) .
\end{equation}
\end{enumerate}
\end{enumerate}

\subsection{Energy Expectation Values}

\begin{enumerate}
\item Any element of $A_{1}\in \mathcal{F}_{1}=$ $\mathcal{F}^{1}$ can be
expanded as 
\begin{equation}
A_{1}=\left( \boldsymbol{x}\right) \mathcal{R}_{x}\left( A_{1}\right)
=\left( 
\begin{array}{ccc}
x_{1} & \cdots & x_{2r}%
\end{array}%
\right) \left( 
\begin{array}{c}
\mathcal{R}_{x}\left( A_{1}\right) _{1} \\ 
\vdots \\ 
\mathcal{R}_{x}\left( A_{1}\right) _{2r}%
\end{array}%
\right)
\end{equation}%
leading to 
\begin{equation}
A_{n}=\left( \boldsymbol{x}_{n}\right) \mathcal{R}_{x}\left( A_{n}\right)
=\left( 
\begin{array}{ccc}
x_{1}\cdots x_{n} & \cdots & x_{2r-n}\cdots x_{2r}%
\end{array}%
\right) \left( 
\begin{array}{c}
\mathcal{R}_{x}\left( A_{n}\right) _{1} \\ 
\vdots \\ 
\mathcal{R}_{x}\left( A_{n}\right) _{\binom{2r}{n}}%
\end{array}%
\right) .
\end{equation}

\item The Clifford complex symmetric inner product 
\begin{equation}
\left( \left. {}\right\vert \right) :\mathcal{F\times F}\rightarrow \mathbb{C%
}
\end{equation}%
is defined by

\begin{enumerate}
\item 
\begin{equation}
\left\langle \left\langle I_{\mathcal{F}}\right\rangle \right\rangle _{0}=1
\end{equation}

\item 
\begin{equation}
\left( \left. a_{1}\right\vert b_{1}\right) =\left\langle \left\langle \frac{%
1}{2}\left\{ a_{1},b_{1}\right\} \right\rangle \right\rangle
_{0}=\left\langle \left\langle \frac{1}{2}\left(
a_{1}b_{1}+b_{1}a_{1}\right) \right\rangle \right\rangle _{0}
\end{equation}

\item 
\begin{equation}
\left( \left. a_{1}a_{2}\cdots a_{n}\right\vert b_{1}b_{2}\cdots
b_{m}\right) =\frac{1}{\sqrt{n!m!}}\det \left\{ \left( 
\begin{array}{ccc}
\left( \left. a_{1}\right\vert b_{1}\right) & \cdots & \left( \left.
a_{1}\right\vert b_{m}\right) \\ 
\vdots & \ddots & \vdots \\ 
\left( \left. a_{n}\right\vert b_{1}\right) & \cdots & \left( \left.
a_{n}\right\vert b_{m}\right)%
\end{array}%
\right) \right\} \delta _{nm}
\end{equation}

\item Let $a,b\in \mathcal{F}^{1}$%
\begin{equation}
\left( \left. a\right\vert b\right) =\mathcal{R}_{x}\left( a\right)
^{t}\left( \left. \mathbf{x}\right\vert \mathbf{x}\right) \mathcal{R}%
_{x}\left( a\right) =\mathcal{R}_{x}\left( a\right) ^{t}\Delta \left( 
\mathbf{x}\right) \mathcal{R}_{x}\left( b\right) ;\quad \mathcal{R}%
_{x}\left( a\right) ,\mathcal{R}_{x}\left( b\right) \in \mathbb{C}^{r},
\end{equation}%
where the metrix matrix $\Delta \left( \mathbf{x}\right) $ is defined as 
\begin{equation}
\Delta \left( \mathbf{x}\right) =\left( \left. \mathbf{x}\right\vert \mathbf{%
x}\right) .
\end{equation}
\end{enumerate}

\item States $\mathcal{S}\left( \mathcal{F}\right) $ associated with $%
\mathcal{F}$ are positive normalized linear functionals $f$ 
\begin{equation}
f\left( A\right) \geq 0\,\forall A\in \mathcal{F}_{+},
\end{equation}%
where the positive cone $\mathcal{F}_{+}$ of $\mathcal{F}$ is defined by 
\begin{equation}
\mathcal{F}_{+}=\left\{ \left. {}\right\vert \right\}
\end{equation}

\item 
\begin{equation}
E_{H}\left( D\right) =\limfunc{Tr}\nolimits_{\mathcal{H}_{\mathcal{F}%
}}\left( HD\right) =\left\langle \left. H\right\vert D\right\rangle _{%
\mathcal{H}_{\mathcal{F}}},\quad D\in S\left( \mathcal{B}\left( \mathcal{H}_{%
\mathcal{F}}\right) \right) ,\quad H=H^{\dagger }\in \mathcal{L}\left( 
\mathcal{H}_{\mathcal{F}}\right)
\end{equation}

\item Let 
\begin{equation}
H=\Gamma \left( h\right)
\end{equation}%
then 
\begin{equation}
E_{H}\left( D\right) =\left\langle \left. \Gamma \left( h\right) \right\vert
D\right\rangle _{\mathcal{H}_{\mathcal{F}}}=\left\langle \left. h\right\vert
L\left( D\right) \right\rangle _{\mathcal{H}_{\mathcal{F}}}.
\end{equation}

\begin{enumerate}
\item The Hamiltonian can be expressed in normal form in terms of \textbf{%
creation and annihilation} operators and the \textbf{operator product }$%
\equiv $ \textbf{Clifford product }as the strongly convergent expansion 
\begin{equation}
H=\Gamma \left( h\right) =\sum_{\substack{ 0\leq q\leq r  \\ 0\leq p\leq r}}%
\sum_{\substack{ 1\leq j_{1}<\cdots <j_{q}\leq r  \\ 1\leq k_{1}<\cdots
<k_{q}\leq r}}a_{j_{1}}^{\dagger }\cdots a_{j_{q}}^{\dagger }a_{k_{p}}\cdots
a_{k_{1}}\mathcal{R}_{a}\left( H\right) _{j_{1}\cdots j_{q}k_{p}\cdots
k_{1}}.
\end{equation}

\item The Hamiltonian can be expressed in normal form in terms of \textbf{%
creation and annihilation} operators and the \textbf{Clifford exterior
product} as the strongly convergent expansion%
\begin{equation}
H=\Gamma \left( h\right) =\sum_{0\leq n\leq 2r}\sum_{\substack{ 1\leq
j_{1}<\cdots <j_{p}\leq r  \\ 1\leq k_{1}<\cdots <k_{q}\leq r  \\ p+q=n}}%
a_{j_{1}}^{\dagger }\wedge \cdots \wedge a_{j_{p}}^{\dagger }\wedge
a_{k_{p}}\cdots \wedge a_{k_{1}}\mathcal{R}_{\wedge a}\left( H\right)
_{j_{1}\cdots j_{p}k_{q}\cdots k_{1}}.
\end{equation}

\item The Hamiltonian can be expressed in terms of \textbf{Majorano}
operators and the \textbf{Clifford product} $\equiv $ \textbf{operator
product }as the strongly convergent expansion%
\begin{equation}
H=\Gamma \left( h\right) =\sum_{0\leq n\leq 2r}\sum_{\substack{ 1\leq
j_{1}<\cdots <j_{n}\leq r}}x_{j_{1}}\wedge \cdots \wedge x_{j_{n}}\mathcal{R}%
_{x}\left( H\right) _{j_{1}\cdots j_{n}}=\sum_{0\leq n\leq 2r}\sum 
_{\substack{ 1\leq j_{1}<\cdots <j_{n}\leq r}}x_{j_{1}}\cdots x_{j_{n}}%
\mathcal{R}_{x}\left( H\right) _{j_{1}\cdots j_{n}}.
\end{equation}
\end{enumerate}
\end{enumerate}

\begin{example}
Let the Hamiltonian $H=H_{\left( 4\right) }=\widehat{P_{\mathcal{A}_{4}}}%
\left( H\right) \in \mathcal{A}_{4}$ then%
\begin{eqnarray}
H &=&\sum_{\substack{ 0\leq q\leq 2  \\ 0\leq p\leq 2}}\sum_{\substack{ %
1\leq j_{1}<\cdots <j_{q}\leq r  \\ 1\leq k_{1}<\cdots <k_{q}\leq r}}%
a_{j_{1}}^{\dagger }\cdots a_{j_{q}}^{\dagger }a_{k_{p}}\cdots a_{k_{1}}%
\mathcal{R}_{a}\left( H\right) _{j_{1}\cdots j_{q}k_{1}\cdots k_{p}} \\
&=&\left\vert \phi \right\rangle \left\langle \phi \right\vert \mathcal{R}%
_{a}\left( H\right) _{00}+\sum_{1\leq j_{1}\leq r}a_{j_{1}}^{\dagger }%
\mathcal{R}_{a}\left( H\right) _{j_{1}0}+\sum_{1\leq j_{1}\leq r}a_{k_{1}}%
\mathcal{R}_{a}\left( H\right) _{0k_{1}} \\
&&+\sum_{1\leq j_{1}<j_{2}\leq r}a_{j_{1}}^{\dagger }a_{j_{2}}^{\dagger }%
\mathcal{R}_{a}\left( H\right) _{j_{1}j_{2}0}+\sum_{\substack{ 1\leq
j_{1}\leq r  \\ 1\leq k_{1}\leq r}}a_{j_{1}}^{\dagger }a_{k_{1}}\mathcal{R}%
_{a}\left( H\right) _{j_{1}k_{1}}+\sum_{1\leq k_{1}<k_{2}\leq
r}a_{k_{1}}^{\dagger }a_{k_{2}}^{\dagger }\mathcal{R}_{a}\left( H\right)
_{0k_{1}k_{2}} \\
&&+\sum_{\substack{ 1\leq j_{1}<j_{2}\leq r  \\ 1\leq k_{1}<\leq r}}%
a_{j_{1}}^{\dagger }a_{j_{2}}^{\dagger }a_{k_{1}}\mathcal{R}_{a}\left(
H\right) _{j_{1}j_{2}k_{1}}+\sum_{\substack{ 1\leq j_{1}\leq r  \\ 1\leq
k_{1}<k_{2}\leq r}}a_{j_{1}}^{\dagger }a_{k_{2}}a_{k_{1}}\mathcal{R}%
_{a}\left( H\right) _{j_{1}k_{1}k_{2}} \\
&&+\sum_{1\leq j_{1}<\cdots <j_{4}\leq r}a_{j_{1}}^{\dagger }\cdots
a_{j_{4}}^{\dagger }\mathcal{R}_{a}\left( H\right) _{j_{1}\cdots j_{4}}+\sum 
_{\substack{ 1\leq j_{1}<\cdots <j_{3}\leq r  \\ 1\leq k_{1}<\leq r}}%
a_{j_{1}}^{\dagger }\cdots a_{j_{3}}^{\dagger }a_{k_{1}}^{\dagger }\mathcal{R%
}_{a}\left( H\right) _{j_{1}\cdots j_{3}k_{1}}+\sum_{\substack{ 1\leq
j_{1}<j_{2}\leq r  \\ 1\leq k_{1}<k_{2}\leq r}}a_{j_{1}}^{\dagger
}a_{j_{2}}^{\dagger }a_{k_{2}}a_{k_{1}}\mathcal{R}_{a}\left( H\right)
_{j_{1}j_{2}k_{1}k_{2}} \\
&&+\sum_{\substack{ 1\leq j_{1}\leq r  \\ 1\leq k_{1}<k_{2}<k_{3}\leq r}}%
a_{j_{1}}^{\dagger }a_{k_{3}}a_{k_{2}}a_{k_{1}}\mathcal{R}_{a}\left(
H\right) _{j_{1}k_{1}k_{2}k_{3}}+\sum_{1\leq k_{1}<k_{2}<k_{3}<k_{4}\leq
r}a_{k_{4}}a_{k_{3}}a_{k_{2}}a_{k_{1}}\mathcal{R}_{a}\left( H\right)
_{k_{1}k_{2}k_{3}k4}.
\end{eqnarray}
\end{example}

Let $H$ be a two particle Hamiltonian $H_{\left( 2\right) }$ then

\begin{enumerate}
\item 
\begin{equation}
H=H_{\left( 2\right) }\in \mathcal{A}_{00}\oplus \mathcal{A}_{11}\oplus 
\mathcal{A}_{22}
\end{equation}%
and can be expressed in

\item wedge products of creation and anninilation operators as 
\begin{eqnarray}
&&H_{\left( 2\right) }=I_{\mathcal{F}}\mathcal{R}_{a}\left( H\right)
_{00}+\sum_{1\leq j_{1}\leq r}a_{j_{1}}^{\dagger }\wedge a_{j_{1}}\mathcal{R}%
_{a}\left( H\right) _{j_{1}j_{1}}+\sum_{1\leq j_{1}\neq k_{1}\leq
r}a_{j_{1}}^{\dagger }\wedge a_{k_{1}}\mathcal{R}_{a}\left( H\right)
_{j_{1}k_{1}} \\
&&+\sum_{\substack{ 1\leq j_{1}<j_{2}\leq r  \\ 1\leq k_{1}<k_{2}\leq r}}%
a_{j_{1}}^{\dagger }\wedge a_{j_{2}}^{\dagger }\wedge a_{k_{2}}\wedge
a_{k_{1}}\mathcal{R}_{a}\left( H\right) _{j_{1}j_{2}k_{2}k_{1}}+\sum 
_{\substack{ 1\leq j_{1}<j_{2}\leq r  \\ 1\leq k_{2}\leq r}}%
a_{j_{1}}^{\dagger }\wedge a_{j_{2}}^{\dagger }\wedge a_{k_{2}}\wedge
a_{j_{1}}\mathcal{R}_{a}\left( H\right) _{j_{1}j_{2}k_{2}j_{1}} \\
&&+\sum_{\substack{ 1\leq j_{1}<j_{2}\leq r  \\ 1\leq k_{1}\leq r}}%
a_{j_{1}}^{\dagger }\wedge a_{j_{2}}^{\dagger }\wedge a_{j_{2}}\wedge
a_{k_{1}}\mathcal{R}_{a}\left( H\right) _{j_{1}j_{2}j_{2}k_{1}}+\sum_{1\leq
j_{1}<j_{2}\leq r}a_{j_{1}}^{\dagger }\wedge a_{j_{2}}^{\dagger }\wedge
a_{j_{2}}\wedge a_{j_{1}}\mathcal{R}_{a}\left( H\right)
_{j_{1}j_{2}j_{2}k_{1}}
\end{eqnarray}%
rearranging one obtains 
\begin{eqnarray}
&&H_{\left( 2\right) }=I_{\mathcal{F}}\mathcal{R}_{a}\left( H\right)
_{00}+\sum_{1\leq j_{1}\leq r}a_{j_{1}}^{\dagger }\wedge a_{j_{1}}\mathcal{R}%
_{a}\left( H\right) _{j_{1}j_{1}}+\sum_{1\leq j_{1}\neq k_{1}\leq
r}a_{j_{1}}^{\dagger }\wedge a_{k_{1}}\mathcal{R}_{a}\left( H\right)
_{j_{1}k_{1}} \\
&&+\sum_{\substack{ 1\leq j_{1}<j_{2}\leq r  \\ 1\leq k_{1}<k_{2}\leq r}}%
a_{j_{1}}^{\dagger }\wedge a_{j_{2}}^{\dagger }\wedge a_{k_{2}}\wedge
a_{k_{1}}\mathcal{R}_{a}\left( H\right) _{j_{1}j_{2}k_{2}k_{1}}+\sum 
_{\substack{ 1\leq j_{1}<j_{2}\leq r  \\ 1\leq k_{2}\leq r}}%
a_{j_{1}}^{\dagger }\wedge a_{j_{1}}\wedge a_{j_{2}}^{\dagger }\wedge
a_{k_{2}}\mathcal{R}_{a}\left( H\right) _{j_{1}j_{2}k_{2}j_{1}} \\
&&+\sum_{\substack{ 1\leq j_{1}<j_{2}\leq r  \\ 1\leq k_{1}\leq r}}%
a_{j_{1}}^{\dagger }\wedge a_{k_{1}}\wedge a_{j_{2}}^{\dagger }\wedge
a_{j_{2}}\mathcal{R}_{a}\left( H\right) _{j_{1}j_{2}j_{2}k_{1}}+\sum_{1\leq
j_{1}<j_{2}\leq r}a_{j_{1}}^{\dagger }\wedge a_{j_{1}}\wedge
a_{j_{2}}^{\dagger }\wedge a_{j_{2}}\mathcal{R}_{a}\left( H\right)
_{j_{1}j_{2}j_{2}k_{1}}
\end{eqnarray}%
after noting 
\begin{equation}
a_{j}^{\dagger }\wedge a_{j}=\left( a_{j}^{\dagger }a_{j}-\frac{1}{2}I_{%
\mathcal{F}}\right)
\end{equation}%
this leads to 
\begin{eqnarray}
&&H_{\left( 2\right) }=I_{\mathcal{F}}\mathcal{R}_{a}\left( H\right)
_{00}+\sum_{1\leq j_{1}\leq r}\left( a_{j_{1}}^{\dagger }a_{j_{1}}-\frac{1}{2%
}I_{\mathcal{F}}\right) \mathcal{R}_{a}\left( H\right)
_{j_{1}j_{1}}+\sum_{1\leq j_{1}\neq k_{1}\leq r}a_{j_{1}}^{\dagger }a_{k_{1}}%
\mathcal{R}_{a}\left( H\right) _{j_{1}k_{1}} \\
&&+\sum_{1\leq j_{1}\neq k_{1}<j_{2}\neq k_{2}\leq r}a_{j_{1}}^{\dagger
}a_{j_{2}}^{\dagger }a_{k_{2}}a_{k_{1}}\mathcal{R}_{a}\left( H\right)
_{j_{1}j_{2}k_{2}k_{1}}+\sum_{1\leq j_{1}\neq k_{2}<j_{2}\neq k_{2}\leq
r}\left( a_{j_{1}}^{\dagger }a_{j_{1}}-\frac{1}{2}I_{\mathcal{F}}\right)
a_{j_{2}}^{\dagger }a_{k_{2}}\mathcal{R}_{a}\left( H\right)
_{j_{1}j_{2}k_{2}j_{1}} \\
&&+\sum_{1\leq j_{1}\neq k_{1}<j_{2}\neq k_{1}\leq r}a_{j_{1}}^{\dagger
}a_{k_{1}}\left( a_{j_{2}}^{\dagger }a_{j_{2}}-\frac{1}{2}I_{\mathcal{F}%
}\right) \mathcal{R}_{a}\left( H\right) _{j_{1}j_{2}j_{2}k_{1}}+\sum_{1\leq
j_{1}<j_{2}\leq r}\left( a_{j_{1}}^{\dagger }a_{j_{1}}-\frac{1}{2}I_{%
\mathcal{F}}\right) \left( a_{j_{2}}^{\dagger }a_{j_{2}}-\frac{1}{2}I_{%
\mathcal{F}}\right) \mathcal{R}_{a}\left( H\right) _{j_{1}j_{2}j_{2}k_{1}}
\end{eqnarray}

\item operator products of creation and anninilation operators as 
\begin{equation}
H_{\left( 2\right) }=\Gamma \left( h_{\left( 2\right) }\right) =\Gamma
\left( \left\vert \phi \right\rangle \left\langle \phi \right\vert \mathcal{R%
}_{a}\left( H\right) _{00}+\sum_{\substack{ 1\leq j_{1}\leq r  \\ 1\leq
k_{1}\leq r}}\left\vert \varphi _{j_{1}}\right\rangle \left\langle \varphi
_{k_{1}}\right\vert \mathcal{R}_{a}\left( H\right) _{j_{1}k_{1}}+\sum 
_{\substack{ 1\leq j_{1}<j_{2}\leq r  \\ 1\leq k_{1}<k_{2}\leq r}}\left\vert
\varphi _{j_{1}}\varphi _{j_{2}}\right\rangle \left\langle \varphi
_{k_{1}}\varphi _{k_{2}}\right\vert \mathcal{R}_{a}\left( H\right)
_{j_{1}j_{2}k_{2}k_{1}}\right) ,
\end{equation}%
where

\item the reduced Hamiltonian $h_{\left( 2\right) }$ is 
\begin{equation}
h_{\left( 2\right) }=\left\vert \phi \right\rangle \left\langle \phi
\right\vert \mathcal{R}_{a}\left( H\right) _{00}+\sum_{\substack{ 1\leq
j_{1}\leq r  \\ 1\leq k_{1}\leq r}}\left\vert \varphi _{j_{1}}\right\rangle
\left\langle \varphi _{k_{1}}\right\vert \mathcal{R}_{a}\left( H\right)
_{j_{1}k_{1}}+\sum_{\substack{ 1\leq j_{1}<j_{2}\leq r  \\ 1\leq
k_{1}<k_{2}\leq r}}\left\vert \varphi _{j_{1}}\varphi _{j_{2}}\right\rangle
\left\langle \varphi _{k_{1}}\varphi _{k_{2}}\right\vert \mathcal{R}%
_{a}\left( H\right) _{j_{1}j_{2}k_{2}k_{1}}
\end{equation}

\item and the state be a $N$ -particle state $D^{N}$ then 
\begin{eqnarray}
E_{H}\left( D^{N}\right) &=&\left\langle \left. \Gamma \left( h_{2}\right)
\right\vert D^{N}\right\rangle _{\mathcal{H}_{\mathcal{F}}}=\left\langle
\left. h_{2}\right\vert L\left( D^{N}\right) \right\rangle _{\mathcal{H}_{%
\mathcal{F}}}=\left\langle \left. h_{2}\right\vert P_{\mathcal{H}%
^{2}}L\left( D^{N}\right) P_{\mathcal{H}^{2}}\right\rangle _{\mathcal{H}_{%
\mathcal{F}}} \\
&=&\left\langle \left. h_{2}\right\vert P_{\mathcal{H}^{2}}L\left(
D^{N}\right) P_{\mathcal{H}^{2}}\right\rangle _{\mathcal{H}_{\mathcal{F}}}=%
\binom{N}{2}\left\langle \left. h_{2}\right\vert L_{N}^{2}\left(
D^{N}\right) \right\rangle _{\mathcal{H}_{\mathcal{F}}}
\end{eqnarray}%
\begin{equation}
D=L^{-1}\left( \Delta \right)
\end{equation}%
\begin{equation}
E_{H}\left( D\right) =\left\langle \left. h\right\vert \Delta \right\rangle
_{\mathcal{H}_{\mathcal{F}}}=\left\langle \left. h\right\vert P_{\mathcal{H}%
^{2}}\Delta P_{\mathcal{H}^{2}}\right\rangle _{\mathcal{H}_{\mathcal{F}}}
\end{equation}

\item 
\begin{equation}
H=\Gamma \left( h\right) =\sum_{0\leq n\leq 2r}\sum_{\substack{ 1\leq
j_{1}<\cdots <j_{n}\leq r}}x_{j_{1}}\wedge \cdots \wedge x_{j_{n}}\mathcal{R}%
_{x}\left( H\right) _{j_{1}\cdots j_{n}}=\sum_{0\leq n\leq 2r}\sum 
_{\substack{ 1\leq j_{1}<\cdots <j_{n}\leq r}}x_{j_{1}}\cdots x_{j_{n}}%
\mathcal{R}_{x}\left( H\right) _{j_{1}\cdots j_{n}}.
\end{equation}%
\begin{equation}
E_{H}\left( D\right) =\mathcal{R}_{x}\left( H\right) _{0}\mathcal{R}%
_{x}\left( D\right) _{0}+\mathcal{\cdots }+\mathcal{R}_{x}\left( H\right)
_{2r}^{t}\mathcal{R}_{x}\left( D\right) _{2r}
\end{equation}%
\begin{equation}
D>0,\limfunc{Tr}D=1=\left\Vert D\right\Vert _{\mathcal{B}_{2}\left( \mathcal{%
H}_{\mathcal{F}}\right) }
\end{equation}%
\begin{equation}
E_{H}\left( D\right) =\mathcal{R}_{x}\left( H\right) _{0}\mathcal{R}%
_{x}\left( D\right) _{0}+\mathcal{R}_{x}\left( H\right) _{1}^{t}\mathcal{R}%
_{x}\left( D\right) _{1}+\mathcal{R}_{x}\left( H\right) _{2}^{t}\mathcal{R}%
_{x}\left( D\right) _{2}
\end{equation}
\end{enumerate}

\subsection{Even Subalgebra}

The even subalgebra $\mathcal{F}_{\text{even}}=C\ell _{2r}^{+}(\mathbb{C})$
consists of elements of even grade. For $r\geq 2$, it is isomorphic to: 
\begin{equation}
\mathcal{F}_{\text{even}}\cong \mathrm{Mat}(2^{r-1},\mathbb{C})\oplus 
\mathrm{Mat}(2^{r-1},\mathbb{C}).
\end{equation}

\subsection{Grade-4 Subspace}

Define the subspace: 
\begin{equation}
\mathcal{F}_{4\text{even}}=\mathcal{F}^{0}\oplus \mathcal{F}^{2}\oplus 
\mathcal{F}^{4}
\end{equation}%
consisting of elements with grades $0,2,$and $4,$ $\mathcal{F}^{0},\mathcal{F%
}^{2},\mathcal{F}^{4}$ respectively.

\section{Positive Cones}

\begin{enumerate}
\item The full positive cones$\mathcal{\ }$of $\mathcal{F}$

\begin{enumerate}
\item The full positive cone of the full algebra $\mathcal{F}_{+}\subset 
\mathcal{F}$ is defined by 
\begin{equation}
\mathcal{F}_{+}=\left\{ \left. x=a^{\dagger }a\right\vert a\in \mathcal{F}%
\right\} .
\end{equation}

\item The full positive cone of the even subalgebra $\mathcal{F}_{\text{even}%
+}\subset \mathcal{F}_{\text{even}}$ is defined by 
\begin{equation}
\mathcal{F}_{\text{even}+}=\left\{ \left. x=a^{\dagger }a\right\vert a\in 
\mathcal{F}_{\text{even}}\right\} .
\end{equation}

\item In general one has the proper inclusion $\mathcal{F}_{\text{even}%
+}\subset \mathcal{F}_{+}$ as

\begin{enumerate}
\item Consider 
\begin{eqnarray}
\left( a_{0}+a_{1}+a_{2}\right) ^{\dagger }\left( a_{0}+a_{1}+a_{2}\right)
&=&a_{0}^{\dagger }a_{0}+a_{0}^{\dagger }a_{1}\ +a_{0}^{\dagger }a_{2}\
+a_{1}^{\dagger }a_{0}+a_{1}^{\dagger }a_{1}\ +a_{1}^{\dagger }a_{2}\
+a_{2}^{\dagger }a_{0}+a_{2}^{\dagger }a_{1}\ +a_{2}^{\dagger }a_{2}\  \\
&=&\left( a_{0}^{\dagger }a_{0}+a_{1}^{\dagger }a_{1}+a_{2}^{\dagger
}a_{2}+a_{0}^{\dagger }a_{2}+a_{2}^{\dagger }a_{0}\right) +\left(
a_{0}^{\dagger }a_{1}+a_{1}^{\dagger }a_{0}+a_{1}^{\dagger
}a_{2}+a_{2}^{\dagger }a_{1}\right)
\end{eqnarray}

\item In general 
\begin{equation}
a=a_{\text{even}}+a_{\text{odd}}
\end{equation}%
giving

\item 
\begin{eqnarray}
\left( a_{\text{even}}+a_{\text{odd}}\right) ^{\dagger }\left( a_{\text{even}%
}+a_{\text{odd}}\right) &=&a_{\text{even}}^{\dagger }a_{\text{even}}+a_{%
\text{odd}}^{\dagger }a_{\text{odd}}+a_{\text{odd}}^{\dagger }a_{\text{even}%
}+a_{\text{even}}^{\dagger }a_{\text{odd}} \\
&=&b=b_{\text{even}}+b_{\text{odd}}\in \mathcal{F}_{+}
\end{eqnarray}%
which gives

\item 
\begin{equation}
b_{\text{even}}=b_{\text{even}}^{\dagger }=a_{\text{even}}^{\dagger }a_{%
\text{even}}+a_{\text{odd}}^{\dagger }a_{\text{odd}}\in \mathcal{F}_{\text{%
even}+}
\end{equation}%
and

\item 
\begin{equation}
b_{\text{odd}}=b_{\text{odd}}^{\dagger }=a_{\text{odd}}^{\dagger }a_{\text{%
even}}+a_{\text{even}}^{\dagger }a_{\text{odd}}
\end{equation}

\item 
\begin{equation}
\limfunc{Tr}\left( b_{\text{odd}}\right) =0
\end{equation}%
thus $b_{\text{odd}}$ must have positive and \textbf{negative} spectrum.

\item Hence in general elements $b$ of $\mathcal{F}_{+}$ contain components $%
b_{\text{odd}}$ that do not belong to $\mathcal{F}_{\text{even}},$ which can
be seen by considering that the sum of two commuting self adjoint operators $%
A$ and $B$ that have spectral decomposions 
\begin{equation}
A=i\sum_{1\leq j<k\leq 2r}c_{jk}x_{j}x_{k}=\int_{a\in \mathbb{R}%
_{+}}aP\left( a\right) da
\end{equation}%
and 
\begin{equation}
B=\sum_{1\leq j\leq 2r}c_{j}x_{j}\ =\int_{b\in \mathbb{R}}bP_{B}\left(
b\right) db,
\end{equation}%
where%
\begin{eqnarray}
a &>&b \\
\left[ P\left( a\right) ,P_{B}\left( b\right) \right] &=&0\quad \\
\forall a,b &\in &\mathbb{R}
\end{eqnarray}%
so that 
\begin{equation}
a+b\geq 0.
\end{equation}
\end{enumerate}
\end{enumerate}

\item The grade 4 \textbf{Reduced Even Cone} $\Pi _{4}\left( \mathcal{F}_{%
\text{even}+}\right) $ (\textbf{Reduced Second Order Positive Cone}) defined
by 
\begin{equation}
\mathcal{C}=\Pi _{4}\left( \mathcal{F}_{\text{even}+}\right) =\mathcal{F}%
_{4}\cap \mathcal{F}_{\text{even}+},
\end{equation}%
where 
\begin{eqnarray}
\Pi _{4} &:&\mathcal{F\rightarrow F}_{4} \\
\Pi _{4}^{2} &=&\Pi _{4}=\Pi _{4}^{\dagger }
\end{eqnarray}

\item Intrinsic cones generated by the subspace $\mathcal{F}_{4}$ are
defined as

\begin{enumerate}
\item 
\begin{equation}
\mathcal{F}_{4+}=\left\{ \left. x=a^{\dagger }a\right\vert a\in \mathcal{F}%
_{4}\right\} \subset \mathcal{F}_{4}
\end{equation}

\item 
\begin{equation}
\mathcal{F}_{2+}=\left\{ \left. x=a^{\dagger }a\right\vert a\in \mathcal{F}%
_{2}\right\} \subset \mathcal{F}_{2}
\end{equation}

\item 
\begin{equation}
\mathcal{F}_{0+}=\left\{ \left. x=a^{\dagger }a\right\vert a\in \mathcal{F}%
_{0}\right\} \subset \mathcal{F}_{0}
\end{equation}

\item 
\begin{equation}
\mathcal{F}_{+}^{4}=\left\{ \left. x=a^{\dagger }a\right\vert a\in \mathcal{F%
}^{4}\right\} \subset \mathcal{F}^{4}
\end{equation}

\item 
\begin{equation}
\mathcal{F}_{+}^{2}=\left\{ \left. x=a^{\dagger }a\right\vert a\in \mathcal{F%
}^{2}\right\} \subset \mathcal{F}^{2}
\end{equation}

\item 
\begin{equation}
\mathcal{F}_{+}^{0}=\left\{ \left. x=a^{\dagger }a\right\vert a\in \mathcal{F%
}^{0}\right\} \subset \mathcal{F}^{0}
\end{equation}
\end{enumerate}

\item Bra-Ket positive cones

\begin{enumerate}
\item 
\begin{enumerate}
\item 
\begin{equation}
\mathcal{S}^{0}=\mathcal{H}^{0}
\end{equation}

\item 
\begin{equation}
\mathcal{S}^{1}=\mathcal{H}^{0}\oplus \mathcal{H}^{1}
\end{equation}

\item 
\begin{equation}
\mathcal{S}^{2}=\mathcal{H}^{0}\oplus \mathcal{H}^{1}\oplus \mathcal{H}^{2}
\end{equation}

\item 
\begin{equation}
\mathcal{B}\left( \mathcal{S}^{n}\right) =\bigoplus\limits_{0\leq p,q\leq n}%
\mathcal{B}\left( \mathcal{H}^{p},\mathcal{H}^{q}\right) ;\quad 0\leq n\leq 2
\end{equation}

\item 
\begin{eqnarray}
P_{S^{0}} &=&P_{H^{0}} \\
P_{S^{1}} &=&P_{H^{0}}+P_{H^{1}} \\
P_{S^{2}} &=&P_{H^{0}}+P_{H^{1}}+P_{H^{2}}
\end{eqnarray}
\end{enumerate}

\item 
\begin{equation}
\mathcal{B}\left( \mathcal{H}^{0}\right) _{+}\oplus \mathcal{B}\left( 
\mathcal{H}^{1}\right) _{+}=P_{\mathcal{H}^{0}\oplus \mathcal{H}^{1}}%
\mathcal{F}_{2+}P_{\mathcal{H}^{0}\oplus \mathcal{H}^{1}}=\left\{ \left.
x=a^{\dagger }a\right\vert a\in \mathcal{F}_{\text{even}}\right\}
\end{equation}

\item 
\begin{equation}
P_{\mathcal{H}^{0}}\mathcal{F}_{0+}P_{\mathcal{H}^{0}}=\left\{ \left.
x=a^{\dagger }a\right\vert a\in \mathcal{F}_{\text{even}}\right\}
\end{equation}

\item 
\begin{equation}
P_{\mathcal{H}^{2}}\mathcal{F}_{+}^{4}P_{\mathcal{H}^{2}}=\left\{ \left.
x=a^{\dagger }a\right\vert a\in \mathcal{F}_{\text{even}}\right\}
\end{equation}

\item 
\begin{equation}
P_{\mathcal{H}^{1}}\mathcal{F}_{+}^{2}P_{\mathcal{H}^{1}}=\left\{ \left.
x=a^{\dagger }a\right\vert a\in \mathcal{F}_{\text{even}}\right\}
\end{equation}

\item 
\begin{equation}
P_{\mathcal{H}^{0}}\mathcal{F}_{+}^{0}P_{\mathcal{H}^{0}}=\left\{ \left.
x=a^{\dagger }a\right\vert a\in \mathcal{F}_{\text{even}}\right\}
\end{equation}
\end{enumerate}

\item Bra-Kets in terms of field operators%
\begin{align}
\left\vert \varphi _{j_{1}}\cdots \varphi _{j_{p}}\right\rangle \left\langle
\varphi _{j_{1}}\cdots \varphi _{j_{p}}\right\vert & =\dprod\limits_{0\leq
l\leq p}n_{i}\dprod\limits_{p+1\leq l\leq r}h_{i}=\dprod\limits_{0\leq l\leq
p}n_{i}\dprod\limits_{p+1\leq l\leq r}\left( I_{\mathcal{F}}-n_{l}\right) \\
& =\dprod\limits_{0\leq l\leq p}\left\{ \frac{1}{2}I_{\mathcal{F}%
}+ix_{l}x_{l+r}\right\} \dprod\limits_{p+1\leq l\leq r}\left\{ \frac{1}{2}I_{%
\mathcal{F}}-ix_{l}x_{l+r}\right\} .
\end{align}
\end{enumerate}

\section{Characterization of Grade-4 Inhomogeneous Projectors}

A self-adjoint projector $P\in \mathcal{F}_{4}$ has the form: 
\begin{equation}
P=a+i\sum_{i<j}b_{ij}x_{i}x_{j}+\sum_{i<j<k<l}c_{ijkl}x_{i}x_{j}x_{k}x_{l}
\end{equation}%
with $a\in \mathbb{R}$, $b_{ij}\in \mathbb{R}$, $c_{ijkl}\in \mathbb{R}$.

\begin{theorem}[Grade 4 Projector Constraints]
The element $P$ in (4) is a projector ($P^{2}=P$) with no components beyond
grade 4 if and only if the following constraints hold: 
\begin{align}
\mathbf{C1(grade0):}& \quad
a^{2}-a=-\sum_{i<j}b_{ij}^{2}+\sum_{i<j<k<l}c_{ijkl}^{2}  \label{eq:C1} \\
\mathbf{C2(grade2):}& \quad
(2a-1)b_{ij}+4\sum_{k<l}\sum_{m}b_{km}c_{mlij}=0\quad \forall i<j
\label{eq:C2} \\
\mathbf{C3(grade4):}& \quad
(2a-1)c_{ijkl}-(b_{ij}b_{kl}-b_{ik}b_{jl}+b_{il}b_{jk})  \notag \\
& \qquad +\langle p_{4}^{2}\rangle _{4}|_{ijkl}+2\langle p_{2}p_{4}\rangle
_{4}|_{ijkl}=0\quad \forall i<j<k<l  \label{eq:C3} \\
\mathbf{C4(grade6):}& \quad \sum_{\text{pairings}}(\pm
)b_{pq}c_{rstu}=0\quad \forall \text{ 6-index sets}  \label{eq:C4} \\
\mathbf{C5(grade8):}& \quad \sum_{\text{4-4 partitions}}(\pm
)c_{I}c_{J}=0\quad \forall \text{ 8-index sets}  \label{eq:C5} \\
\mathbf{Rank:}& \quad 2^{2r}a\in \{0,1,2,\ldots ,2^{2r}\}  \label{eq:rank}
\end{align}
\end{theorem}

\begin{remark}
Condition \eqref{eq:C5} is precisely the quadratic Pl\`{u}cker relations
characterizing that $p_{4}=\sum c_{ijkl}x_{i}x_{j}x_{k}x_{l}$ corresponds to
a point on the Grassmannian $\mathrm{Gr}(4,2r)$ of 4-planes in $\mathbb{C}%
^{2r}$.
\end{remark}

\begin{remark}
A self-adjoint projector $P\in \mathcal{F}_{even}$ has the form: 
\begin{gather}
P=c_{0}I_{\mathcal{F}}+i\sum_{1\leq j_{1}<j_{2}\leq
2r}c_{j_{1}j_{2}}x_{j_{1}}x_{j_{2}}+\cdots +\sum_{1\leq j_{1}<\cdots
<j_{2r}\leq 2r}c_{j_{1}\cdots j_{2r}}x_{j_{1}}\cdots x_{j_{2r}} \\
=P_{0}+\cdots +P_{2r}
\end{gather}%
with $c_{0}\in \mathbb{R}$,$c_{j_{1}\cdots }c_{j_{2p}}\in \mathbb{R}$.
\end{remark}

\begin{remark}
\begin{eqnarray}
P &=&P^{2} \\
P_{0}+\cdots +P_{2p} &=&\left( P_{0}+\cdots +P_{2r}\right) ^{2}=\sum_{0\leq
p\leq r}P_{2p}\left( P_{0}+\cdots +P_{2r}\right) =\sum_{0\leq p\leq
r}\left\{ P_{2p}P_{0}+\cdots +P_{2p}P_{2r}\right\} =\sum_{0\leq p,q\leq
r}P_{2p}P_{2q}
\end{eqnarray}%
\begin{equation}
P_{2p}P_{2q}=\left\langle P_{2p}P_{2q}\right\rangle _{\left\vert
2p-2q\right\vert }+\left\langle P_{2p}P_{2q}\right\rangle _{\left\vert
2p-2q\right\vert +2}+\left\langle P_{2p}P_{2q}\right\rangle _{\left\vert
2p-2q\right\vert +4}+\left\langle P_{2p}P_{2q}\right\rangle _{\left\vert
2p-2q\right\vert +6}+\left\langle P_{2p}P_{2q}\right\rangle _{\left\vert
2p-2q\right\vert +8}+\cdots +\left\langle P_{2p}P_{2q}\right\rangle _{2p+2q}
\end{equation}
\end{remark}

\begin{remark}
\begin{equation}
P_{2p}=\sum_{\substack{ 0\leq q\leq r  \\ 0\leq s\leq r}}\Theta \left(
q+s-p\right) \Theta \left( p+\left\vert q-s\right\vert \right) \left\langle
P_{2q}P_{2s}\right\rangle _{2p}
\end{equation}
\end{remark}

\begin{remark}
\begin{equation}
P_{2p}P_{2q}=\sum_{\left\vert p-q\right\vert \leq m\leq p+q}\left\langle
P_{2p}P_{2q}\right\rangle _{2m}
\end{equation}
\end{remark}

\begin{remark}
\begin{equation}
p+q+\left\vert p-q\right\vert =2\max (p,q)
\end{equation}%
\begin{equation}
p+q-\left\vert p-q\right\vert =2\min \U{2061} (p,q)
\end{equation}%
\begin{equation}
P_{2p}=\sum_{0\leq q\leq r}\left\langle P_{2p}P_{2q}\right\rangle _{2\left(
2\max (p,q)-q-\left\vert p-q\right\vert \right) }
\end{equation}%
\begin{equation}
c_{0}=\sum_{0\leq p\leq r}\sum_{1\leq j_{1}<\cdots <j_{2p}\leq
2p}c_{j_{1}\cdots j_{2p}}^{2}
\end{equation}
\end{remark}

\begin{remark}
\begin{equation}
c_{j_{1}\cdots j_{2p}}=\sum_{\substack{ 0\leq q\leq r  \\ 0\leq s\leq r}}%
\sum _{\substack{ \text{contractions }  \\ k=q+s-p}}\text{sign}\left(
a_{1}\cdots a_{2q}b_{1}\cdots b_{2q}\right) c_{a_{1}\cdots
a_{2q}}c_{b_{1}\cdots b_{2s}}
\end{equation}
\end{remark}

\begin{theorem}[Reduced Projector Constraints]
The element $P$ in $\mathcal{F}_{even}$ is a projector ($P^{2}=P$)
\end{theorem}

\section{Partial States on $\mathcal{F}_{4}$}

A partial normal state $\omega $ on $\mathcal{F}_{4}$ is given by a density
matrix: 
\begin{equation}
d=\frac{1}{2^{2r}}I+i\sum_{i<j}\beta _{ij}x_{i}x_{j}+\sum_{i<j<k<l}\delta
_{ijkl}x_{i}x_{j}x_{k}x_{l}
\end{equation}%
with $\beta _{ij}\in \mathbb{R}$, $\delta _{ijkl}\in \mathbb{R}$, and $%
\mathrm{Tr}(d)=1$ (automatically satisfied by the normalization).

\section{The Extension Problem}

Let $\mathcal{P}_{4}$ denote the set of all self-adjoint projectors in $%
\mathcal{F}_{4}$ satisfying \eqref{eq:C1}-\eqref{eq:rank}.

\begin{theorem}[Krein Extension Condition]
The partial state $\omega $ with density matrix $d$ extends to a state on
the full even subalgebra $\mathcal{F}$ if and only if: 
\begin{equation}
\max_{P\in \mathcal{P}_{4}}\mathrm{Tr}(dP)\leq 1
\end{equation}
\end{theorem}

\begin{theorem}[Explicit Trace Formula]
For $P \in \mathcal{P}_4$ and $d$ as in (6): 
\begin{equation}
\mathrm{Tr}(dP) = a - 2^{2r}\sum_{i<j}\beta_{ij}b_{ij} -
2^{2r}\sum_{i<j<k<l}\delta_{ijkl}c_{ijkl}
\end{equation}
\end{theorem}

\begin{remark}
Thus the extension condition becomes: 
\begin{equation}
\max_{(a,b_{ij},c_{ijkl})\in \mathcal{P}_{4}}\left( a-2^{2r}\sum_{i<j}\beta
_{ij}b_{ij}-2^{2r}\sum_{i<j<k<l}\delta _{ijkl}c_{ijkl}\right) \leq 1
\end{equation}
\end{remark}

\section{The Number Operator}

Define the number operator in the fermionic Fock representation: 
\begin{equation}
N = \sum_{k=1}^{r} \psi_k^\dagger \psi_k
\end{equation}

In terms of Clifford generators: 
\begin{equation}
N=\frac{r}{2}I-\frac{i}{2}\sum_{k=1}^{r}x_{2k-1}x_{2k}\in \mathcal{F}_{4}
\end{equation}

\begin{theorem}[Integrality Condition]
\begin{equation}
\mathrm{Tr}(Nd) = \frac{r}{2} - 2^{2r-1}\sum_{k=1}^{r} \beta_{2k-1,2k} = n
\in \mathbb{Z}_{>0}
\end{equation}
\end{theorem}

\section{Geometric Interpretation}

\begin{theorem}[Interior Characterization]
If $\mathrm{Tr}(Nd)=n\in \mathbb{Z}_{>0}$ with $0<n<2r$, then $d$ lies in
the interior of the state space when restricted to $\mathcal{F}_{4}$.
\end{theorem}

\begin{proof}
For $r>2$, primitive idempotents (pure states) contain components of grade $%
>4$, hence do not lie in $\mathcal{F}_{4}$. Therefore no pure state
restricts to an extremal state on $\mathcal{F}_{4}$. The integrality
condition forces $d$ to be a mixture of Fock states (eigenstates of $N$),
which have integer expectation values. Such mixtures lie in the interior of
the convex set of states on $\mathcal{F}_{4}$.
\end{proof}

\section{Canonical Extensions}

Since $d$ is in the interior, many extensions exist. A canonical choice is
the grand canonical ensemble:

\begin{equation}
D_\beta = \frac{e^{-\beta N}}{\mathrm{Tr}(e^{-\beta N})}
\end{equation}

where $\beta \in \mathbb{R}$ is chosen to satisfy: 
\begin{equation}
\mathrm{Tr}(N D_\beta) = n
\end{equation}

This gives the equation: 
\begin{equation}
\frac{\sum_{k=0}^{2r} k e^{-\beta k} \binom{2r}{k}}{\sum_{k=0}^{2r}
e^{-\beta k} \binom{2r}{k}} = n
\end{equation}

\section{Summary}

\begin{tabular}{|l|l|}
\hline
\textbf{Concept} & \textbf{Mathematical Formulation} \\ \hline
Partial state & $d=2^{-2r}I+d_{2}+d_{4}$ \\ 
Projector in $\mathcal{F}_{4}$ & $P=a+p_{2}+p_{4}$ with constraints C1-C5 \\ 
Extension condition & $\max_{P\in \mathcal{P}_{4}}(a-2^{2r}\langle
d_{2},p_{2}\rangle -2^{2r}\langle d_{4},p_{4}\rangle )\leq 1$ \\ 
Number operator & $N=\frac{r}{2}I-\frac{i}{2}\sum_{k}x_{2k-1}x_{2k}$ \\ 
Integrality & $\frac{r}{2}-2^{2r-1}\sum_{k}\beta _{2k-1,2k}=n\in \mathbb{Z}%
_{>0}$ \\ 
Consequence & $d$ in interior \^{a}\ddag ' many extensions \\ 
Canonical extension & $D_{\beta }=e^{-\beta N}/\mathrm{Tr}(e^{-\beta N})$ \\ 
\hline
\end{tabular}

\section*{Acknowledgments}

This work emerged from extensive mathematical discussions on the structure
of Clifford algebras and the extension theory of states on operator algebras.

\begin{thebibliography}{99}
\bibitem{Krein} Krein, M. G. (1940). Sur le probl\`{e}me du prolongement des
fonctions hermitiennes positives et continues. Comptes Rendus (Doklady) de
l'Acad\'{e}mie des Sciences de l'URSS, 26, 17--22.

\bibitem{Clifford} P. Lounesto, \textit{Clifford Algebras and Spinors},
Cambridge University Press, 2001.

\bibitem{Fock} V. Fock, \textquotedblleft Konfigurationsraum und zweite
Quantelung,\textquotedblright\ \textit{Zeitschrift f\~{A}%
%TCIMACRO{\U{bc}}%
%BeginExpansion
$\frac14$%
%EndExpansion
r Physik}, 1932.

\bibitem{} @book\{Naimark1972,

\bibitem{} author = \{M. A. Naimark\},

\bibitem{} title = \{Normed Algebras\},

\bibitem{} publisher = \{Wolters-Noordhoff\},

\bibitem{} year = \{1972\},

\bibitem{} note = \{Section 30, Theorem 1\}

\bibitem{} \}

\bibitem{} @book\{Paulsen2002,

\bibitem{} author = \{Vern I. Paulsen\},

\bibitem{} title = \{Completely Bounded Maps and Operator Algebras\},

\bibitem{} publisher = \{Cambridge University Press\},

\bibitem{} series = \{Cambridge Studies in Advanced Mathematics\},

\bibitem{} volume = \{78\},

\bibitem{} year = \{2002\},

\bibitem{} chapter = \{3\}

\bibitem{} \}

\bibitem{} @book\{Blackadar2006,

\bibitem{} author = \{Bruce Blackadar\},

\bibitem{} title = \{Operator Algebras: Theory of \{\$C\symbol{94}%
*\$\}-Algebras and von Neumann Algebras\},

\bibitem{} publisher = \{Springer\},

\bibitem{} series = \{Encyclopaedia of Mathematical Sciences\},

\bibitem{} volume = \{122\},

\bibitem{} year = \{2006\},

\bibitem{} note = \{Section II.6.3\}

\bibitem{} \}

\bibitem{} @incollection\{Krein1983,

\bibitem{} author = \{M. G. Krein\},

\bibitem{} title = \{On the Problem of the Extension of a Hermitian-Positive
Continuous Function\},

\bibitem{} booktitle = \{Selected Works, Volume 3, Part 2\},

\bibitem{} publisher = \{Gordon and Breach\},

\bibitem{} year = \{1983\},

\bibitem{} note = \{Original paper: C. R. (Doklady) Acad. Sci. URSS 26
(1940), 17--22\}

\bibitem{} \}

\bibitem{} @article\{ErdahlJin2000,

\bibitem{} author = \{R. M. Erdahl and B. Jin\},

\bibitem{} title = \{The lower bound method for reduced density matrices\},

\bibitem{} journal = \{Journal of Mathematical Physics\},

\bibitem{} volume = \{41\},

\bibitem{} number = \{6\},

\bibitem{} pages = \{3857--3879\},

\bibitem{} year = \{2000\}

\bibitem{} \}

\bibitem{} @article\{Mazziotti2012,

\bibitem{} author = \{David A. Mazziotti\},

\bibitem{} title = \{Structure of fermionic density matrices: complete
\{\$N\$\}-representability conditions\},

\bibitem{} journal = \{Physical Review Letters\},

\bibitem{} volume = \{108\},

\bibitem{} number = \{26\},

\bibitem{} pages = \{263002\},

\bibitem{} year = \{2012\}

\bibitem{} \}

The first four entries are the core theoretical references; the last two are
examples of their application to the N\U{2011}representability problem,
included in case they are useful.
\end{thebibliography}

\end{document}
